\subsection{Prose Chaining (WikiText-103)}
\label{sec:prose_chaining}

\paragraph{Task Description.}
The prose chaining experiment evaluates deep multi-step generation on
\texttt{wikitext-103-raw-v1}, a large Wikipedia-derived corpus with markedly
broader and more diverse vocabulary than wikitext-2.  The increased lexical
distribution creates a denser but flatter value attractor field, making
chaining harder: the system must bridge further in attractor space to
reconstruct the held-out 60\% suffix of each sample.

wikitext-103 was chosen specifically because its long-tail vocabulary
represents the regime where shallow n-gram statistics break down but
structural value resonance remains viable --- making it a sharper
discriminator for the architecture's generative capabilities.

\paragraph{Results.}
Across $N = 10$ test samples the mean weighted score was 0.266.

\paragraph{Assessment.}
Partial chaining was observed.  Common Wikipedia structural patterns
(section openers, reference sentence rhythms, passive constructions) are
recoverable, but long-range semantic coherence is limited by the attractor
density achievable with 10 ingested samples from a 103-million-token corpus.

Figure~\ref{fig:prose_chaining_map} shows the trial outcome map.
